\documentclass[a4paper,twoside]{article}


% --- RUIMTE EN LAYOUT ---
\usepackage{setspace}
\setstretch{1.15}

\usepackage{geometry} %size regulating
\geometry{
 a4paper,
 left=25mm,
 right=25mm,
 top=20mm,
 bottom=20mm,
 }

\usepackage{parskip} % for better spacing between paragraphs


% --- WISKUNDE ENSYMBOLEN ---
\usepackage{amsmath, amssymb, amsfonts, amsthm} % For nicer math format
\usepackage{mathrsfs} % For curly arrr
\usepackage{bm} % For bold symbols eg \bm{\nabla}
\usepackage{esint} % For bolletje in integrals
\usepackage{dsfont}
\usepackage{siunitx} % Voeg deze expliciet toe voor die graden en eenheden!


% --- OVERIGE ---

\usepackage{graphicx} % Required for inserting images
\usepackage[hidelinks]{hyperref} % geen kaders om linken
\usepackage{subcaption}
%\usepackage[dutch]{babel}
\usepackage[T1]{fontenc}
\usepackage[utf8]{inputenc}
\usepackage{lmodern}
\usepackage{multicol}

% --- Tikz packages ---
\usepackage{tikz}
\usetikzlibrary{angles,quotes, babel} % for pic (angle labels)
\usetikzlibrary{arrows.meta,decorations.markings,calc, decorations.pathmorphing, decorations.pathreplacing, positioning, shapes, shadings, patterns}
\usepackage{xcolor}
\usepackage{tikz-3dplot}
\usepackage[siunitx]{circuitikz}

\usepackage{etoolbox} % ifthen
\usepackage[outline]{contour} % glow around text
\usepackage{xfp} % higher precision (16 digits?)
\contourlength{1.1pt}



% --- THEOREM STIJLEN MET EXTRA RUIMTE ---

\theoremstyle{definition}
\newtheorem{postulate}{Postulaat}
\newtheorem{definition}{Definitie}
\newtheorem{derivation}{Afleiding}

\theoremstyle{plain}
\newtheorem{theorem}{Theorem}


% --- EXTRA COMMANDO'S ---

\newcommand{\dbar}{{\mkern+3mu\mathchar'26\mkern-12mu d}}

% --- Titel en Introductie --- 

\title{Formularium voor oefeningen Sterrenstelsels}
\author{T. Cools}
\date{ 2025-2026}

\begin{document}
\maketitle

\begin{abstract}

Deze notities bevatten alle formules die van belang zijn om de oefeningen van Sterrenstelsels op te lossen.
Deze worden dan ook gedacht te kennen te zijn als ze niet in het formulrium staan, maar er wordt geen garantie gegeven dat dit ook echt het geval is.

1 van de redenen dat ik dit maak is omdat niet alle formules in de cursus worden gegeven.
\end{abstract}

\section{Oefenset 1}

Schijnbare magnitude:
\begin{equation}
    m_X (d) = -2.5 \log \left( \frac{F_X (d)}{F_{X,0}} \right)
\end{equation}

Absolute magnitude:
\begin{equation}
    M_X = m_X(10 \text{ pc}) = -2.5 \log \left( \frac{F_X (10 \text{ pc})}{F_{X,0}} \right)
\end{equation}

\begin{equation}
    M_X - m_X = - 5 \log \left( \frac{d}{10 \text{ pc}} \right)
\end{equation}  

Merk op dat:
\begin{align*}
    M_X - m_X &= - 5 \log \left( \frac{d [pc]}{10 [pc]} \right) \\
    &= - 5 \log \left( \frac{d}{10} \right) \\
    &= - 5 \left[\log \left( d \right) - 1\right]  \\
\end{align*}


Herinner van sterren en planeten dat:
\begin{align*}
    dL &= -L\alpha(r)dr\\
    \implies \ln \left( \frac{L}{L_0} \right) &= -\int_0^R \alpha(r) dr\\
    \implies L &= L_0 e^{-\int_0^R \alpha(r) dr} = L_0 e^{-\tau}\\
    &\boxed{L = L_0 e^{-\tau}}
\end{align*}
Waarbij $\tau$ de optische diepte is. Soms ook \(\tau_X\) om de band aan te duiden.


FLux-Lichtkracht relatie:
\begin{equation}
    F_X = \frac{L_X}{4 \pi d^2}
\end{equation}

En:
\[
F_X(\tau_X) = \frac{L_X(\tau_X)}{4 \pi d^2} = \frac{L_{X} e^{-\tau_X}}{4 \pi d^2}
\]
Dus:
\[
m_X (\tau_X) = -2.5 \log \left( \frac{F_X  e^{-\tau_X}}{F_{X,0}} \right) = m_X + 2.5 \log \left( e^{-\tau_X} \right)
\]


\section{Oefenset 2}

De centrale logaritmische oppervlaktehelderheid \(\mu_X (R=0)\):
\begin{equation}
    \mu_X (0) = m_X + 2.5 \log \left(\Omega \right)
\end{equation}

Of soms:
\[ 
\mu_X = -2.5 \log \left(\frac{S_X}{F_{X,0}}\right)
\]

De lineaire oppervlaktehelderheid \(S_X [mag/arcsec^2]\):
\begin{equation}
    S_X = \frac{F_X}{\Omega} = \frac{L_X}{4 \pi d^2 \Omega} = \frac{L_X}{4 \pi^2 r_e^2}
\end{equation}

Relatie tussen \(\mu_X\) en \(S_X\):
\begin{equation}
    \mu_X = -2.5 \log \left( \frac{S_X}{F_{X,0}} \right) = -2.5 \log \left( \frac{F_X/\Omega}{F_{X,0}} \right)   
\end{equation}  


\[
\Omega[arcsec^2] = \Omega [sterad] \left(\frac{3600 \cdot 180}{\pi}\right)^2  
\]

Dit is handig omdat:
\[
d^2 \Omega [sterad] = A
\]
Waarbij \(A\) de oppervlakte is van het object. En dus:
\[
\Omega [sterad] = \frac{A}{d^2}
\]

De intensiteit \(I_X\):
\begin{equation}
    I_X = \frac{L_X}{A} = \frac{L_X}{d^2 \Omega [sterad]}
\end{equation}

Schaallengte \(h\):
\begin{align*}
    \int_0^{L_X} d L_X &= \int_0^{R_{max}} I_X (R) 2 \pi R dR\\
    & = 2\pi I_X(0) \int_0^{R_{max}} e^{-R/h}  R dR \\
    & = 2\pi I_X(0) h^2 \int_0^{\frac{R_{max}}{h}} \exp \left( -\frac{R}{h} \right) \left( \frac{R}{h} \right) d\left( \frac{R}{h} \right)\\
\end{align*}

De \(R_{25}\)-isofoot is de straal waarbij de oppervlaktehelderheid 25 mag/arcsec\(^2\) is. Dus:
\begin{equation}
    \mu_X (R_{25}) = 25 \text{ mag/arcsec}^2
\end{equation}

2 handige vergelijkingen die niet in het formularium staan:

\[
L_\lambda = \int_{0}^{\infty} I(R) 2 \pi R dR
\]
en 
\[
L(R_e) = \frac{1}{2} L_\lambda
\]

\section{Oefenset 3}



Interstellair stralingsveld \(J_\nu\):
\begin{equation}
    J_\nu 4 \pi = \int I_\nu d\Omega
\end{equation}  

\[
4 \pi J_{em} = \int I_{em} d\Omega = \int I_{em} \frac{dA}{r^2} = \int \frac{L_{em}}{4 \pi r^2} \frac{dA}{r^2} = \int \frac{L_{em}}{4 \pi r^4} dA
\]


Deze hadden we random nodig:

\[
\int_{0}^{\infty}\kappa_\lambda B_\lambda d \lambda 
\]


\section{Oefenset 4}

\[
\kappa_\lambda = \kappa_0 \left(\frac{\lambda_0}{\lambda}\right)^\beta
\]
met \(\beta = 2\) voor de emissie en \(\beta = 0\) voor de absorptie term.

Conditie voor thermisch evenwicht:
\begin{equation}
    J^\text{abs} = J^\text{em}
\end{equation}  

\section{Oefenset 5}

Straal HIII-regio:
\begin{equation}
    R_s = \left( \frac{3 Q_0}{4 \pi n_H^2 \alpha_B (T)} \right)^{1/3}
\end{equation}  
met \(Q_0\) het aantal ioniserende fotonen per seconde, \(n_H\) de waterstofdichtheid en \(\alpha_B\) de recombinatiecoëfficiënt.

Bijbehorende massa:
\begin{equation}
    M_s = \frac{4}{3} \pi R_S^3 n_H m_H = \frac{m_H Q_0}{n_H \alpha_B (T)}
\end{equation}

\section{Oefenset 6}

Weten dat de golflengte van de H\(\alpha\)-lijn 656 nm is en weten dat als we werken met warm geïoniseerd gas, dat de temperatuur daar dan ongeveer \(10^4\) K is.

Als de Gaussische verdeling van de fluxdichtheid \(F_\nu\) er zo uitziet:
\[
F_\nu = \frac{A}{\sqrt{2 \pi} \sigma \nu_0} \exp \left( -\frac{(\nu - \nu_0)^2}{2 \sigma^2 \nu_0^2} \right)
\]
dan is de totale flux \(F\):
\[
F = \int_{-\infty}^{\infty} F_\nu d\nu = A
\]
Wat een hele hoop rekenwerk bespaart (tenzij ze expliciet vraagt om \(F\) te berekenen).

Opnieuw twee belangrijke vergelijkingen die niet in het formularium staan, maar die we wel nodig hadden:
\[
\langle v \rangle = \int_0^\infty  F (v) v dv
\]
en 
\[
\sigma^2 = \int_{-\infty}^{+\infty} F(v) (v - \langle v \rangle)^2 dv
\]

We hebben ook 
\[
\int_{-\infty}^{+\infty} F(v) dv = 1
\]
gebruikt, want \(F(v)\) is een genormaliseerde kansverdeling.

\[
\int_{-\infty}^{+\infty} F(v) v^2 dv = \langle v \rangle^2
\]

Dit zijn blijkbaar allemaal definities van de cursus statistiek (dat nog niet iedereen heeft gevolgd).
\end{document}